\chapter{INTRODUCTION}
\label{chap:intro}




Since the start of scientific thinking, mathematics has been one of the key languages in which we can understand the world, and Partial Differential Equations (PDEs) are one of the main dialects of the mathematical language we use to model and to simplify nature's complexity.

When modelling these equations, we cannot always afford them to have simple closed-form solutions (that is, when they have a solution at all). This exposes the need to obtain approximate solutions coming from numerical methods for these PDEs, also giving the option to circumvent analytical complexity and allowing the extraction of predictions, comparisons with data, the building of simulations, and so on. These classical numerical methods, widely studied for all types of differential equations, usually rely on discretizing the PDE's domain into grids and performing temporal steps.

With the recent ascension of Machine Learning (ML) methods and a whole ecosystem of tools built around it, a natural path is to bridge these two areas and let ML augment the numerical solution of PDEs, taking advantage of the many software frameworks developed. This is one facet of Scientific Machine Learning (SciML), the broader effort to fuse physics-based, mechanistic models with data-driven methods.

Within SciML, this work focuses on the numerical solution of the Boussinesq System. Given the vast literature around the application of SciML on PDEs, we restrict our literature review to works on hydrodynamics equations.

\section{Literature Overview on Scientific Machine Learning for Hydrodynamics PDEs}

The neural network building blocks underlying all of these approaches trace a line beginning with the formal neuron of \textcite{mcculloch1943} and Rosenblatt's perceptron \parencite{rosenblatt1958}, the first trainable single-layer architecture. The key algorithmic advance that unlocked multi-layer training was the backpropagation algorithm \parencite{rumelhart1986}, making gradient-based optimization practical for deep architectures. Shortly after, the universal approximation theorem \parencite{cybenko1989, hornik1989} provided the theoretical guarantee for a multilayer perceptron (MLP) with a single hidden layer of sufficient width to be able to approximate any continuous function on a compact domain to arbitrary precision. This combination of practical trainability and universal approximation established the MLP as the backbone of modern deep learning. The term neural network itself reflects the biological metaphor at the heart of these architectures, in which \textcite{mcculloch1943} drew an explicit analogy to the firing of biological neurons, and this vocabulary propagated through the field, causing MLPs to be routinely and interchangeably referred to as artificial neural networks (ANNs) or feedforward neural networks. The idea of ``deep'' in deep learning, popularized following the breakthrough results of \textcite{lecun2015deep}, simply denotes an MLP with many hidden layers, a regime in which representational capacity grows substantially and where the backpropagation-based training paradigm proves to be exceptionally efficient.

These ML foundations become the basis of a scientific methodology once the methods are coupled to the physical laws that govern a problem of interest. This is the premise of Scientific Machine Learning, a term consolidated by \textcite{baker2019} to name the discipline at the intersection of scientific computing and machine learning, and surveyed in breadth by \textcite{karniadakis2021}. Its unifying goal is to combine the interpretability, generalization, and inductive biases of physics-based models with the flexibility and data-adaptivity of ML, rather than treating the network as a pure black box. The methods span several families. Some recover governing equations directly from data, as in the Sparse Identification of Nonlinear Dynamics (SINDy) \parencite{brunton2016}. Others embed differential-equation structure into the network itself, as in Neural Ordinary Differential Equations (neural ODEs) \parencite{chen2018node} and Universal Differential Equations (UDEs), in which a trainable network stands in for the unknown terms of an otherwise mechanistic model \parencite{rackauckas2020}. These tools have been carried across a wide range of domains, from global weather and climate forecasting \parencite{bonev2023} to molecular dynamics, materials design, and turbulence modelling \parencite{karniadakis2021}. Among all these directions, the one that concerns us here is the use of SciML to obtain numerical solutions of hydrodynamics equations, and in particular of the Boussinesq system presented in Section~\ref{sec:intro_boussinesq}.

SciML for PDEs was reshaped by one idea in particular, the Physics-Informed Neural Networks (PINNs) \parencite{raissi2019}, that is, the embedding of the governing PDE residual directly into the training loss of a Neural Network, working as a regularization term for not only fitting data, but also respecting the underlying physics. In hydrodynamics, a growing body of work has applied PINNs to water flow \parencite{depina2021, depina2023}, shallow-water dynamics on the sphere \parencite{bihlo2022}, and inverse problems in strongly nonlinear wave systems \parencite{jagtap2022}. The consistent result that we can address from reviews \parencite{cuomo2023} is that PINNs are most effective when observational data are sparse and the task is parameter identification, namely inverse problems, which are usually not well-posed. But, for a reliable precise forward simulation at scale, they do not yet rival classical solvers. In our vision, these approaches serve best to complement them.

As a natural extension from neural networks, there is a recent development in learning maps between infinite-dimensional function spaces, in which the theoretical foundation was popularized by \textcite{chen1995}, who extended the universal approximation theorem to nonlinear operators, in other words, a shallow network can approximate any continuous operator between Banach spaces to arbitrary precision. Bringing this idea into practical terms, \textcite{lu2021deeponet} introduced DeepONet, the first deep architecture expressly designed for learning operators coming from differential equations, demonstrating it across a range of ODEs and PDEs. Building on this line of work, operator learning puts aside the single-instance limitation of PINNs by training an operator to approximate the full mapping from parameters and initial data to the solution, so that at inference time a new solution costs only a forward pass for the network. Parallel to DeepONet, Fourier Neural Operator (FNO) \parencite{li2021} is one of the most widely adopted instances of this idea: it parametrizes an integral kernel in Fourier space, achieving speed-ups over other neural operators, and it is discretization-independent, the so-called zero-shot superresolution. Subsequent works explored alternative bases like wavelets \parencite{gupta2021}, but FNO was validated on many types of equations such as dispersive and soliton wave equations \parencite{pei2022, zhong2023}, and was also scaled to geophysical problems spanning regional ocean modeling, shallow-water emulation, and global weather forecasting \parencite{bonev2023, lkhagvasuren2024, uchida2025, patel2025, yu2025}.

Following the idea for PINNs, the Physics-Informed Neural Operator (PINO) \parencite{li2023pino} combines operator learning with PDE-residual constraints, cutting the labelled-data requirement while preserving physical consistency. Applications to shallow-water systems \parencite{rosofsky2023} confirmed that this hybrid regime inherits the generalization of FNO without its data hunger. Broad surveys of the field \parencite{piml2023review} identify these physics-data approaches working together as a promising SciML direction, while naming turbulence modeling, scalability, and noisy-data robustness as open challenges.

A known limitation that cuts across all SciML architectures is the so-called spectral bias, the tendency of gradient-based optimizers to learn the low-frequency components of a solution faster than the high-frequency ones \parencite{rahaman2019, xu2019frequency}. The theoretical account of this behavior came with the Neural Tangent Kernel (NTK) \parencite{jacot2018}, which showed that a sufficiently wide network trains, to leading order, as a linear model governed by a kernel that stays essentially fixed during optimization, the so-called lazy-training regime. In this picture each mode of the solution is learned at a rate set by its associated NTK eigenvalue, so the wide disparity between the largest and smallest eigenvalues translates directly into the frequency imbalance seen during training \parencite{arora2019, cao2019, bordelon2020}. \textcite{wang2022ntk} carried this analysis to PINNs, showing that the eigenvalue disparity between the PDE-residual and initial-condition loss components is what drives their spectral bias. This pathology is especially consequential for dispersive wave problems, where the solution energy is broadly distributed across wavenumbers. A recent systematic study \parencite{khodakarami2026} showed that second-order optimizers combined with spectrum-aware loss formulations substantially mitigate it.

As far as this review reveals, SciML methods have been applied to a broad class of shallow-water and dispersive wave equations. To the best of our knowledge, however, no study has systematically evaluated FNO, PINO, or PINNs on the parametric 1D Boussinesq system, where both the nonlinearity coefficient $\alpha$ and the dispersion coefficient $\beta$ vary simultaneously. This gap motivates our effort to implement and evaluate these SciML methods on that system, documenting the details a reader would need to apply them to a dispersive wave problem of their own.

\section{The Boussinesq System of Equations}
\label{sec:intro_boussinesq}

For this work, we choose to implement SciML methods for a specific PDE originating from J.~Boussinesq's 1872 effort to give a theoretical account of the solitary wave of translation observed by J.~Scott Russell in a canal in 1834 \parencite{boussinesq1872}. Russell had noticed that a hump of water, set in motion by a stopping barge, could travel for miles without changing shape, a phenomenon the linear wave theory of the time could not explain. Boussinesq showed that the stability of such a wave is the result of a precise balance between nonlinear steepening and frequency dispersion, two effects that are separately destabilizing but mutually compensating in the shallow-water, long-wave regime. The mathematical structure of this balance, and of dispersive wave theory in general, are treated in depth by \textcite{whitham1974} and \textcite{debnath2012}.

A modern two-field form of the system, for the free-surface elevation $\eta(x,t)$ and the depth-averaged horizontal velocity $u(x,t)$, was systematically derived and classified by \textcite{bona2002}, who identified a four-parameter family of equivalent Boussinesq systems obtainable from the incompressible, irrotational Euler equations by asymptotic simplifications. Earlier, \textcite{peregrine1967} had already written down the standard form of these equations for water of variable depth, establishing the notation and scaling that remain in widespread use. What makes the Boussinesq system particularly suitable as a parametric testbed for this work is that changing the nonlinearity and dispersion coefficients continuously deforms the solution from the near-linear, wave-packet regime to the strongly nonlinear soliton regime, giving rise to a more interesting range of behaviors within a single mathematical framework.

The specific one-dimensional form of the system used in this work, with its parameterization of the nonlinearity coefficient $\alpha$ and the dispersion coefficient $\beta$, was adopted from \textcite{martins2023, martins2024}, which is:
\begin{equation}\label{eq:boussinesq_intro}
\begin{cases}
    \eta_t + u_x + \alpha\,(\eta u)_x = 0, \\[4pt]
    u_t - \dfrac{\beta}{3}\,u_{xxt} + \eta_x + \alpha\,u u_x = 0. \\
\end{cases}
\end{equation}

For this work, the parameters $\alpha$ and $\beta$ in~\eqref{eq:boussinesq_intro} are held equal, defining a one-parameter family of problems whose difficulty grows with the parameter. Namely, for small $\alpha = \beta$ the dynamics are nearly linear and the solution is well-captured by low-frequency modes, while for large $\alpha = \beta$ the strong coupling between nonlinearity and dispersion produces richer dynamics that are fundamentally harder to learn. This controlled variation is precisely what makes the setting useful for comparing SciML architectures, with the goal of observing how each method learns the target solution's frequency.

Classical numerical methods for the Boussinesq system rely on pseudospectral and finite-difference discretizations that exploit the periodic or nearly periodic spatial structure of the waves. Pseudospectral \parencite{trefethen2000} schemes achieve spectral accuracy on the spatial derivatives and are therefore well suited to the dispersive character of the equations \parencite{engsigkarup2012}; this is the approach used to generate the reference dataset in this work. On the application side, operational coastal-engineering codes such as FUNWAVE \parencite{funwave2010} implement Boussinesq-type solvers for harbor and nearshore wave propagation, demonstrating the practical relevance of the model beyond the academic setting.

SciML methods have recently been applied only to the Boussinesq equation, not to the specific Boussinesq system. \textcite{zheng2023} applied parameter-integrated PINNs with phase-domain decomposition to reconstruct two-dimensional Boussinesq dynamics including phase transitions and rogue-wave formation, identifying spectral bias as a key obstacle at high nonlinearity. \textcite{wang2024} trained a physics-informed network to infer velocity and pressure fields from surface-elevation data alone on a Boussinesq model, demonstrating the inverse-problem capabilities of the approach. Complementary studies on related nonlinear wave equations \parencite{zhang2024, kumar2025} have confirmed that purely data-driven and purely physics-driven networks each display characteristic failure modes at strong nonlinearity. Taken together, these results suggest that a systematic, parametric evaluation of FNO/PINO/PINN on the 1D Boussinesq system, which is the focus of this work, can be fruitful to understand better the interactions between the methods' ease of training and nonlinearities and dispersiveness.

\section{Research Goals}

This work implements and evaluates four SciML architectures on the parametric 1D Boussinesq system, arranged into two families of three trained methods each. The neural networks family represents a single parameter: a plain multilayer perceptron (MLP) trained on data alone, and a Physics-Informed Neural Network (PINN) trained both with and without supervised data. The operator family learns the whole parameter range in a single training and answers a new parameter with one forward pass: a Fourier Neural Operator (FNO) trained on data alone, and a Physics-Informed Neural Operator (PINO) trained both with and without supervised data. Crossing the two families with the three training regimes, pure data, data and physics, and pure physics, gives the six methods compared throughout. A separate strand runs beside them: a minimal network fitting a fixed time Boussinesq elevation profile, whose Neural Tangent Kernel (NTK) is measured directly to isolate the spectral-bias mechanism. All results are measured against a pseudospectral reference solver, which sets a high-accuracy baseline. The study is guided by the following questions.

\begin{itemize}
    \item Accuracy relative to the pseudospectral baseline: how closely can each family reproduce the reference solutions as nonlinearity and dispersion grow? Where does each method start to fail, and how does its error scale with the difficulty of the regime? For the operators, does the discretization invariance their spectral layers promise survive a query off the training grid?

    \item Data-driven, physics-informed, and hybrid training: how does the presence or absence of supervised reference data affect the accuracy and training behavior of each method? Does adding a physics residual term improve generalization, or does it introduce new difficulties?

    \item Spectral bias across training regimes: how does spectral bias manifest in each architecture when the error spectrum is tracked band by band during training, and to what extent does the inclusion of data or physics constraints alter the frequency at which errors concentrate?

    \item The mechanism behind spectral bias: does Neural Tangent Kernel (NTK) theory genuinely predict the frequency imbalance seen on the Boussinesq system, or only describe it after the fact? We put the theory to two quantitative tests on a minimal network fitting a single Boussinesq elevation profile at a fixed time, across a range of network widths. The first targets the initial-kernel (lazy-training) assumption the analysis depends on: we measure how far the network parameters and the empirical kernel drift over training, and whether the kernel's eigenvalue spectrum is essentially unchanged from initialization to final iteration. The second targets the predicted driver of the bias: modes carrying the high-frequency detail are learned far more slowly than low-frequency ones. Testing both, where they hold and where they break, ties the frequency-band errors seen during training to a concrete mechanism instead of a loose analogy.
\end{itemize}

Together, these questions locate where each family and each training regime holds up on dispersive, nonlinear wave problems and where it fails, and what that implies for choosing among them.

\section{Work Structure}

This work is written with the goal of introducing SciML to senior undergraduates in mathematics, physics, or engineering who aim to apply these methods in their future research, presenting its pros and cons. We assume the reader has good knowledge of scientific computing and linear algebra.
The remaining chapters are organized as follows.

\begin{itemize}
    \item Chapter~2 develops the mathematical background. It covers the
          Boussinesq system and its analytical properties, the discrete Fourier transform and the pseudospectral solver built on it, and then the theoretical foundations of the two families, the MLP and the PINN on one side and the FNO and the PINO on the other. It closes with the Neural Tangent Kernel and a derivation of the spectral bias mechanism that the later chapters put to the test.

    \item Chapter~3 describes the experimental methodology. It details the computational setup, the pseudospectral dataset that serves as both training target and ground truth, the error metrics that score every method on one scale, the six trained methods on the two-family by three-regime grid, and the minimal experiment that isolates the spectral-bias mechanism.

    \item Chapter~4 presents the numerical results. It accounts for what each method costs to train, then takes them along one axis at a time: the nonlinearity tests, the multi-resolution tests, and the spectrum compared mode by mode and closes with the NTK measurements and the frequency-band tracking that tie the observed errors to spectral bias for all methods.

    \item Chapter~5 assesses each family against the questions posed above, states what the study cannot claim, and outlines directions for future work.
\end{itemize}